\documentclass[11pt,a4paper]{article}
\usepackage[margin=22mm]{geometry}
\usepackage{lmodern}
\usepackage{longtable}
\usepackage{hyperref}
\title{B1 LoRMA+ statute classifier: Methodology and Evidence}
\author{CS781 Capstone -- Krishna K Karthikeya G N}
\date{12 November 2026}
\begin{document}
\maketitle
\noindent\textbf{Status.} Post-run methodology reconstruction grounded in archived source and reports. This document is not an executed historical notebook.\par
\section*{Research question}
Test whether a multiplicative low-rank adaptation changes the inherited B1 representation more effectively than additive LoRA. Only the adaptation mechanism changes relative to B1 LoRA; that makes the pair useful for testing parameter-efficient adaptation rather than changing the task pipeline all at once.
\section*{Architecture and mechanism}
LoRMA+ wraps a linear output with trainable low-rank multiplicative factors while preserving the pretrained Phase 1 weights. The same multi-label calibration and evaluation route is retained for a meaningful comparison.
\section*{Measured result}
\begin{center}\begin{tabular}{ll}Scored submission & 1016282\\
Codabench macro-F1 & 0.2542\\
Archived local metric & 0.2517 official-dev macro-F1\end{tabular}\end{center}
The observed Codabench score, 0.2542, is slightly below LoRA's 0.2554. The difference is small and cannot establish a general ranking of the methods.
\section*{Limitations and next check}
A low-rank output correction cannot recover evidence that the window selection never passes to the encoder. Local, pilot, and hidden metrics are deliberately kept distinct. A single scored run does not establish a causal architecture effect or statistical significance.
\section*{Source and reproducibility}
Original source modules: phase5\_b1\_train.py, phase5\_lorma.py, phase5\_core.py. SHA-256 hashes, the original report, and the prediction ZIP receipt are in the companion provenance file. The explanatory notebook was written after the historical run; original execution output is never implied. Reproduction requires authorized course data, the Phase 1 parent weights/tokenizer, and the dependencies in the archived source.
\section*{Documentation disclosure}
AI assistance was used to organize the archived evidence and draft this explanatory document. All numerical claims above come from saved run records or the provided finished-submission screenshot; missing evidence is labelled as such.
\end{document}
